% Einheit 1 von 3 – Irrationale Zahlen und Zahlbereiche (bank/reelle-zahlen/e1.jsonl, zusammenbau v0.1)
%% TODO Zweigzeile Teil 1: was hier gelernt wird (Satz in Schülersprache aus dem Einheitstitel) – Einheit 1
\einheitenkopf[e1]{Einheit 1 von 3 · Irrationale Zahlen und Zahlbereiche}
\zweigzeile{ab Kl. 9 (am Gymnasium ab Kl. 8) · keine P10-Aufgabe}
\setcounter{aufgabe}{9}

%% TODO Titel: Ich-kann-Satz fehlt in der Bank (Platzhalter: Kettenname) – Nr. 10
%% TODO Anweisung: ein Satz über der ersten Teilaufgabe fehlt in der Bank – Nr. 10
\begin{aufgabe}{Zahlbereiche und irrationale Zahlen}
%% TODO \rechenplatz{n}: Schrittzahl des Grundfalls fehlt in der Bank (nur bei fester Schreibform, 2.3 b)
\begin{teile}
\teil $\sqrt{121}$ – geht die Wurzel auf? \\ \kreuz{geht auf} \\ \kreuz{geht nicht auf}
\teil $\frac{5}{8}$ – als Dezimalzahl, abbrechend oder periodisch? Schreibe periodische Zahlen mit Periodenstrich. \\ \kreuz{abbrechend} \\ \kreuz{periodisch} \leerfeld
\teil $\sqrt{169}$ – rational oder irrational? \\ \kreuz{rational} \\ \kreuz{irrational}
\teil $\sqrt{5}$ mit dem Taschenrechner – auf drei Stellen gerundet, mit $=$ oder $\approx$? \leerfeld
\teil $-8$ – in welche Zahlbereiche? Kreuze in der Tabelle an.

\sachtabelle{ccccc}{Zahl & ℕ & ℤ & ℚ & ℝ}{$-8$ & & & & }
\teil $\sqrt{6}$ – markiere die Zahl auf dem Zahlenstrahl. Zwischen welchen Zehnteln liegt sie?

\zahlenstrahl[xmin=2,xmax=3,xstep=0.1,karo=1]{}
\leerfeld $< \sqrt{6} $<$$ \leerfeld
\teil $\sqrt{8}$; $2{,}8$; $2\frac{5}{6}$ – aufsteigend geordnet? \leerfeld $<$ \leerfeld $<$ \leerfeld
\teil Ein Quadrat hat den Flächeninhalt $7\,\text{cm}^2$. Seitenlänge exakt und auf zwei Stellen gerundet? $a =$ \leerfeld $\approx$ \leerfeld
\teil $\sqrt{7} \cdot \sqrt{7}$ – exaktes Ergebnis? Rational oder irrational? \leerfeld
\teil Einschachtelung von $\sqrt{2}$: Fülle die Tabelle aus. Zwischen welchen Zehnteln liegt $\sqrt{2}$?

\sachtabelle{cc}{$x$ & $x^2$}{$1{,}4$ & \\ $1{,}5$ & }
\leerfeld $< \sqrt{2} $<$$ \leerfeld
\teil $\sqrt{40}$: Ist die Zahl rational oder irrational? Begründe. Gib einen Näherungswert auf zwei Stellen an. Zu welchen Zahlbereichen gehört sie?
\end{teile}
\end{aufgabe}

%% TODO Titel: Ich-kann-Satz fehlt in der Bank (Platzhalter: Kettenname) – Nr. 11
%% TODO Anweisung: ein Satz über der ersten Teilaufgabe fehlt in der Bank – Nr. 11
\begin{aufgabe}{Dezimalzahl als Bruch schreiben (rational)}
\begin{teile}
\teil $0{,}45$ – als gekürzter Bruch? \leerfeld
\end{teile}
\end{aufgabe}

%% TODO Titel: Ich-kann-Satz fehlt in der Bank (Platzhalter: Kettenname) – Nr. 12
%% TODO Anweisung: ein Satz über der ersten Teilaufgabe fehlt in der Bank – Nr. 12
\begin{aufgabe}{Teilmengenkette ℕ ⊂ ℤ ⊂ ℚ ⊂ ℝ erklären (jede natürliche Zahl ist ganz, jede ganze rational, jede rationale reell)}
\begin{teile}
\teil Ist jede ganze Zahl rational? Begründe.
\end{teile}
\end{aufgabe}

%% TODO Titel: Ich-kann-Satz fehlt in der Bank (Platzhalter: Kettenname) – Nr. 13
%% TODO Anweisung: ein Satz über der ersten Teilaufgabe fehlt in der Bank – Nr. 13
\begin{aufgabe}{Zahlbereiche und irrationale Zahlen – Fehler finden, Begründen, Darstellungswechsel, Anwendung}
\begin{teile}
\teil Lena schreibt $\sqrt{3} = 1{,}732050808$, weil der Taschenrechner das anzeigt. Finde den Fehler und schreibe richtig.
\teil Warum kann $\sqrt{2}$ kein Bruch sein? Erkläre die Idee in Worten.
\teil Am Zahlenstrahl ist der Punkt A markiert. Welche Zahl ist A? \\ \kreuz{$\sqrt{5}$} \\ \kreuz{$\sqrt{8}$} \\ \kreuz{$\sqrt{11}$}

\zahlenstrahl[xmin=0,xmax=4,xstep=0.5,karo=1]{2.83/A}
\teil Ein quadratischer Teppich soll genau $6\,\text{m}^2$ groß sein. Wie lang muss eine Seite sein? Gib den exakten Wert und einen Wert für den Zuschnitt an.
\end{teile}
\end{aufgabe}
